\thispagestyle{plain}
{\centering
    \vspace*{0.1cm}
    {\Huge \textbf{\color{teal} \coursenamecl} \par}
    {\huge \textbf{{\exercise} 04} \par}
    \vspace{0.4cm}
    {\LARGE \textcolor{ForestGreen}{\studentname} \par}
    % {\LARGE Lớp: \class. \LARGE \faculty \par}
    % {\LARGE \university \par}
    \vspace{0.3cm}
    {\LARGE \today \par}
    \vspace{0.8cm}
}

Em đã thực hiện đầy đủ các yêu cầu của bài và đã làm thêm một số nội dung bổ sung. Em xin chia thành các phần bên dưới để trình bày cụ thể hơn: phần thực hiện yêu cầu Câu 1 (Data Layer, Repository, Service, Transaction), phần thực hiện các Feature của Câu 2 (gồm 3 Feature bắt buộc và 1 Feature khuyến khích), và phần trả lời các câu hỏi Problem Solving (Câu 3).

\section{THỰC HIỆN CÁC YÊU CẦU CỦA BÀI}

\subsection{Link GitHub Repository}

\textbf{Mã nguồn:} \href{https://github.com/qanh-khtn/mini-course-catalog-mvc}{GitHub Repo} (nhánh \texttt{lab04-mini-course}).

\subsection{Tổng quan nâng cấp từ Lab03 lên Lab04}

Ở Lab03, ứng dụng \textbf{Mini Training Center Course Catalog MVC} là một website MVC có Layout dùng chung, Partial View, Tag Helpers, form GET/POST, Model Binding và Validation. Tuy nhiên toàn bộ dữ liệu khóa học vẫn được lưu bằng \texttt{List<Course>} hard-code trong Service, nên dữ liệu bị mất khi tắt ứng dụng.

Sang Lab04, em tiếp tục phát triển trên chính project Lab03 và nâng cấp ứng dụng từ dạng lưu dữ liệu trong bộ nhớ thành một ứng dụng có \textbf{Data Layer rõ ràng, dùng database thật (SQLite) và có luồng xử lý nghiệp vụ an toàn}. Các nội dung nâng cấp chính bao gồm:

\begin{itemize}[leftmargin = 0.6cm, label = -]
    \item Tổ chức lại thư mục thành các tầng \texttt{Models / Data / Repositories / Services / Options / ViewModels}.
    \item Xây dựng các Entity và khai báo Relationship bằng EF Core.
    \item Tạo \texttt{AppDbContext}, cấu hình Mapping và Connection String.
    \item Áp dụng \textbf{Options Pattern} để đọc cấu hình strongly-typed thay vì string key.
    \item Đăng ký \texttt{DbContext}, Repository và Service trong DI Container với lifetime phù hợp.
    \item Tạo Migration, Update Database và Seed Data.
    \item Sử dụng \texttt{AsNoTracking} cho truy vấn chỉ đọc và Tracking cho nghiệp vụ cập nhật.
    \item Xây dựng nghiệp vụ đăng ký khóa học có \textbf{Transaction}, Commit và Rollback.
    \item Tách dependency để thể hiện Minimal Test Mindset (có thể thay bằng Fake Repository).
\end{itemize}

\subsection{Tổ chức thư mục theo tầng Data Layer}

Em đã tổ chức lại project theo đúng yêu cầu của Lab04. Mỗi tầng đảm nhận một trách nhiệm riêng:

\begin{itemize}[leftmargin = 0.6cm, label = -]
    \item \texttt{Models}: chứa các Entity \texttt{Course}, \texttt{CourseCategory}, \texttt{Student}, \texttt{Enrollment}.
    \item \texttt{Data}: chứa \texttt{AppDbContext} và cấu hình Mapping, Seed Data.
    \item \texttt{Repositories}: chứa interface và implementation cho từng Repository.
    \item \texttt{Services}: chứa nghiệp vụ; Controller chỉ gọi Service, không tự query database.
    \item \texttt{Options}: chứa class cấu hình strongly-typed \texttt{TrainingCenterConfig}.
    \item \texttt{ViewModels}: chứa các ViewModel đóng gói dữ liệu trả về View.
\end{itemize}

\begin{figure}[H]
    \centering
    \includegraphics[width=0.9\textwidth]{img/Lab04/folder.jpg}
    \caption{Cấu trúc thư mục project sau khi tổ chức thành các tầng Models, Data, Repositories, Services, Options, ViewModels}
\end{figure}

\subsection{Tạo branch mới cho Lab04}

Theo yêu cầu, em tạo một branch riêng để thực hiện Lab04 nhằm tách biệt với code Lab03 trên nhánh \texttt{main}.

\begin{figure}[H]
    \centering
    \includegraphics[width=0.9\textwidth]{img/Lab04/git_checkout.jpg}
    \caption{Tạo và chuyển sang branch \texttt{lab04-mini-course} để thực hiện Lab04}
\end{figure}

\subsection{Cài đặt EF Core và các package cần thiết}

Em cài đặt các package EF Core để xây dựng Data Layer: \texttt{Microsoft.EntityFrameworkCore}, \\ \texttt{Microsoft.EntityFrameworkCore.Sqlite}, \texttt{Microsoft.EntityFrameworkCore.Design} và công cụ dòng lệnh \texttt{dotnet-ef} để tạo và chạy Migration.

\begin{figure}[H]
    \centering
    \includegraphics[width=0.9\textwidth]{img/Lab04/ef_core_1.jpg}
    \caption{Cài đặt package \texttt{Microsoft.EntityFrameworkCore}}
\end{figure}

\begin{figure}[H]
    \centering
    \includegraphics[width=0.9\textwidth]{img/Lab04/ef_core_3.jpg}
    \caption{Cài đặt package \texttt{Microsoft.EntityFrameworkCore.Design} phục vụ Migration}
\end{figure}

\begin{figure}[H]
    \centering
    \includegraphics[width=0.9\textwidth]{img/Lab04/ef_core_4.jpg}
    \caption{Cài đặt công cụ \texttt{dotnet-ef} để tạo và cập nhật Migration}
\end{figure}

\subsection{Xây dựng Entity và khai báo Relationship}

Em xây dựng 4 Entity có ý nghĩa nghiệp vụ và khai báo các quan hệ giữa chúng:

\begin{itemize}[leftmargin = 0.6cm, label = -]
    \item \texttt{CourseCategory} (chuyên ngành) quan hệ \textbf{One-to-Many} với \texttt{Course}: một chuyên ngành có nhiều khóa học.
    \item \texttt{Course} (khóa học) quan hệ \textbf{One-to-Many} với \texttt{Enrollment}: một khóa học có nhiều lượt đăng ký.
    \item \texttt{Student} (học viên) quan hệ \textbf{One-to-Many} với \texttt{Enrollment}: một học viên có nhiều lượt đăng ký.
\end{itemize}

Như vậy \texttt{Enrollment} là bảng trung gian thể hiện quan hệ \textbf{Many-to-Many} giữa \texttt{Course} và \texttt{Student} thông qua hai quan hệ One-to-Many. Ứng dụng có 1 \texttt{AppDbContext}, 4 Entity và 3 Relationship, đáp ứng yêu cầu tối thiểu của Lab04.

\subsection{Tạo AppDbContext, DbSet và cấu hình Mapping}

Em tạo \texttt{AppDbContext} kế thừa \texttt{DbContext}, khai báo các \texttt{DbSet} cho 4 Entity và cấu hình Mapping bằng Fluent API trong \texttt{OnModelCreating}: khai báo khóa ngoại, hành vi xóa (\texttt{OnDelete}) và kiểu dữ liệu \texttt{decimal(18,2)} cho \texttt{TuitionFee}. Connection String được cấu hình trong \texttt{appsettings.json} với khóa \texttt{DefaultConnection} trỏ tới file SQLite \texttt{MiniCourseCatalog.db}.

\subsection{Áp dụng Options Pattern}

Em tạo class \texttt{TrainingCenterConfig} trong thư mục \texttt{Options} để đọc cấu hình dưới dạng strongly-typed thay vì dùng string key rải rác. Cấu hình gồm \texttt{LowSeatThreshold} (ngưỡng cảnh báo sắp hết chỗ) và \texttt{CenterName}. Class này được bind từ section \texttt{TrainingCenterConfig} trong \texttt{appsettings.json} và được inject vào \texttt{CourseService} qua \texttt{IOptions<TrainingCenterConfig>}.

\subsection{Đăng ký DI Container với Lifetime phù hợp}

Trong \texttt{Program.cs}, em đăng ký các thành phần với lifetime phù hợp:

\begin{itemize}[leftmargin = 0.6cm, label = -]
    \item \texttt{AppDbContext}: \textbf{Scoped} (mặc định của \texttt{AddDbContext}), mỗi HTTP request dùng chung một instance.
    \item Các Repository (\texttt{ICourseRepository}, \texttt{IEnrollmentRepository}, \texttt{IStudentRepository}, \\ \texttt{ICourseCategoryRepository}): \textbf{Scoped}.
    \item Các Service (\texttt{ICourseService}, \texttt{IEnrollmentService}, \texttt{IStudentService}): \textbf{Scoped}.
\end{itemize}

Việc dùng Scoped giúp Service và Repository trong cùng một request chia sẻ chung một \texttt{DbContext}, đảm bảo tính nhất quán của dữ liệu và đúng phạm vi transaction.

\subsection{Repository và Service}

Mỗi Repository có interface và implementation riêng, được inject qua constructor. Repository chịu trách nhiệm truy vấn EF Core; Service chịu trách nhiệm xử lý nghiệp vụ và map sang ViewModel. Controller chỉ gọi Service, không tự tạo Repository bằng \texttt{new} và không query \texttt{DbContext} trực tiếp. Luồng xử lý tuân theo: \textbf{Controller $\rightarrow$ Service $\rightarrow$ Repository $\rightarrow$ DbContext $\rightarrow$ Database}.

\subsection{Migration, Update Database và Seed Data}

Em tạo các Migration có tên rõ nghĩa và chạy Update Database để sinh ra schema thật trong SQLite. Seed Data được khai báo bằng \texttt{HasData} trong \texttt{OnModelCreating}, gồm 4 chuyên ngành, 4 khóa học và 15 học viên, đủ để kiểm tra danh sách, relationship và transaction.

\begin{figure}[H]
    \centering
    \includegraphics[width=0.9\textwidth]{img/Lab04/dotnet_ef_miggration.jpg}
    \caption{Tạo Migration ban đầu bằng lệnh \texttt{dotnet ef migrations add}}
\end{figure}

\begin{figure}[H]
    \centering
    \includegraphics[width=0.9\textwidth]{img/Lab04/migration_update_db_1.jpg}
    \caption{Chạy \texttt{dotnet ef database update} để áp dụng Migration vào database SQLite}
\end{figure}

\begin{figure}[H]
    \centering
    \includegraphics[width=0.9\textwidth]{img/Lab04/seed_data_1.jpg}
    \caption{Dữ liệu Seed (chuyên ngành, khóa học, học viên) đã được ghi vào database}
\end{figure}

\begin{figure}[H]
    \centering
    \includegraphics[width=0.9\textwidth]{img/Lab04/seed_data_2.jpg}
    \caption{Kiểm tra dữ liệu mẫu trong các bảng của database sau khi Seed}
\end{figure}

\subsection{Tracking và AsNoTracking}

Em sử dụng \texttt{AsNoTracking} cho các truy vấn chỉ đọc như danh sách khóa học, trang chi tiết, trang tìm kiếm và trang lịch sử đăng ký, giúp EF Core không cần theo dõi thay đổi nên nhẹ và nhanh hơn. Ngược lại, ở nghiệp vụ đăng ký khóa học em dùng truy vấn \textbf{Tracking} (không có \texttt{AsNoTracking}) để EF Core theo dõi entity \texttt{Course} và tự sinh câu lệnh \texttt{UPDATE} khi giảm số chỗ trống.

\subsection{Transaction cho nghiệp vụ đăng ký khóa học}

Nghiệp vụ trọng tâm của Lab04 là \textbf{đăng ký khóa học}, gồm nhiều bước ghi dữ liệu. Trong \\ \texttt{EnrollmentService.EnrollStudentAsync}, em bọc toàn bộ trong một transaction bằng \\ \texttt{BeginTransactionAsync}:

\begin{enumerate}[leftmargin = 0.7cm, label = \textbf{Bước \arabic*:}]
    \item Truy vấn (Tracking) khóa học cần đăng ký; nếu không tồn tại thì Rollback.
    \item Kiểm tra học viên tồn tại; nếu không thì Rollback.
    \item Kiểm tra khóa học còn chỗ (\texttt{CurrentEnrollment < MaxCapacity}); nếu đầy thì Rollback.
    \item Kiểm tra học viên chưa đăng ký trùng; nếu trùng thì Rollback.
    \item Thêm bản ghi \texttt{Enrollment} mới và tăng \texttt{CurrentEnrollment} của khóa học.
    \item \texttt{SaveChangesAsync} rồi \texttt{CommitAsync}. Nếu có lỗi ở bất kỳ bước nào, \texttt{catch} sẽ gọi \texttt{RollbackAsync} để hủy toàn bộ thay đổi.
\end{enumerate}

Nhờ transaction, nghiệp vụ đảm bảo tính toàn vẹn: hoặc cả hai thao tác (thêm \texttt{Enrollment} và giảm số chỗ) cùng thành công, hoặc cùng bị hủy, không bao giờ rơi vào trạng thái nửa đúng nửa sai.

\subsection{Data Health}

Em xây dựng trang \texttt{/DataHealth} để kiểm tra tình trạng dữ liệu: trạng thái Migration, số lượng bản ghi đã Seed, và minh họa sự khác nhau giữa truy vấn Tracking và AsNoTracking.

\subsection{Minimal Test Mindset}

Toàn bộ Service đều nhận dependency qua constructor dưới dạng interface (\texttt{IEnrollmentRepository}, \\ \texttt{ICourseRepository}\dots) nên có thể thay implementation thật bằng Fake Repository khi kiểm tra logic. Đây là nền tảng để em phát triển thành project unit test thật ở phần làm thêm.

\subsection{Kết quả chạy ứng dụng}

Sau khi hoàn thành các yêu cầu, em chạy ứng dụng bằng \texttt{dotnet run}. Hệ thống khởi tạo server tại \texttt{http://localhost:5063}.

\begin{figure}[H]
    \centering
    \includegraphics[width=0.9\textwidth]{img/Lab04/dotnet_run.jpg}
    \caption{Chạy ứng dụng bằng \texttt{dotnet run}, server khởi tạo thành công}
\end{figure}

\begin{figure}[H]
    \centering
    \includegraphics[width=0.9\textwidth]{img/Lab04/web_home.jpg}
    \caption{Endpoint \texttt{/}: Trang chủ hệ thống Mini Training Center}
\end{figure}

\begin{figure}[H]
    \centering
    \includegraphics[width=0.9\textwidth]{img/Lab04/web_courses.jpg}
    \caption{Endpoint \texttt{/Courses}: Trang danh sách khóa học lấy dữ liệu thật từ database (truy vấn AsNoTracking)}
\end{figure}

\begin{figure}[H]
    \centering
    \includegraphics[width=0.9\textwidth]{img/Lab04/web_courses_detail1.jpg}
    \caption{Endpoint \texttt{/Courses/Detail}: Trang chi tiết hiển thị dữ liệu liên quan qua Relationship (khóa học và chuyên ngành)}
\end{figure}

\begin{figure}[H]
    \centering
    \includegraphics[width=0.9\textwidth]{img/Lab04/web_course_categories.jpg}
    \caption{Endpoint \texttt{/CourseCategories}: Trang thể hiện quan hệ One-to-Many giữa CourseCategory và Course}
\end{figure}

\begin{figure}[H]
    \centering
    \includegraphics[width=0.9\textwidth]{img/Lab04/web_courses_enroll_1.jpg}
    \caption{Endpoint \texttt{/Courses/Enroll}: Form đăng ký khóa học, thực hiện nghiệp vụ nhiều bước có Transaction}
\end{figure}

\begin{figure}[H]
    \centering
    \includegraphics[width=0.9\textwidth]{img/Lab04/web_courses_enroll_2_2.jpg}
    \caption{Đăng ký \textbf{thành công}: Transaction Commit, thêm Enrollment và giảm số chỗ trống của khóa học}
\end{figure}

\begin{figure}[H]
    \centering
    \includegraphics[width=0.9\textwidth]{img/Lab04/web_courses_enroll_3_2.jpg}
    \caption{Trường hợp \textbf{Rollback}: khóa học PRG-201 đã đầy chỗ (30/30), hệ thống từ chối đăng ký và hủy toàn bộ thay đổi}
\end{figure}

\begin{figure}[H]
    \centering
    \includegraphics[width=0.9\textwidth]{img/Lab04/web_data_health.jpg}
    \caption{Endpoint \texttt{/DataHealth}: Trang kiểm tra Migration, Seed Data, tracking/no-tracking và transaction}
\end{figure}

\subsection{Đẩy code lên GitHub}

Em commit và đẩy code lên GitHub với các commit message rõ nghĩa, mỗi nội dung lớn là một commit riêng.

\begin{figure}[H]
    \centering
    \includegraphics[width=0.9\textwidth]{img/Lab04/github_1.jpg}
\end{figure}

\begin{figure}[H]
    \centering
    \includegraphics[width=0.9\textwidth]{img/Lab04/github_4.jpg}
\end{figure}

\section{THỰC HIỆN CÁC TÍNH NĂNG FEATURE (CÂU 2)}

Câu 2 yêu cầu thực hiện 3 Feature bắt buộc và 1 Feature khuyến khích. Ngoài ra em cũng bổ sung thêm một số tính năng đúng trọng tâm Lab04 (Unit Test, Chart.js, UI hiện đại). Mỗi tính năng được commit riêng và đều đi đúng luồng \textbf{Controller $\rightarrow$ Service $\rightarrow$ Repository $\rightarrow$ DbContext}.

\begin{enumerate}[label = \textbf{\arabic*.}, leftmargin = 0.5cm]
    \item \textbf{Project Unit Test thật với xUnit và Fake Repository (Tính năng bổ sung)}
        \begin{itemize}[label = $\bullet$, leftmargin = 0.3cm]
            \item \textbf{Tạo project test riêng:} Em tạo project \texttt{MiniCourseCatalog.Tests} dùng xUnit, tham chiếu tới project chính để kiểm tra logic Service.

            \item \textbf{Viết các Fake Repository:} Em viết \texttt{FakeCourseRepository}, \\ \texttt{FakeCourseCategoryRepository} lưu dữ liệu in-memory, và \\ \texttt{ThrowingEnrollmentRepository} cố tình ném exception ở bước ghi dữ liệu. Việc thay được Repository thật bằng Fake chứng minh kiến trúc đã tách dependency đúng cách.

            \item \textbf{Kiểm thử theo Arrange - Act - Assert:} Em viết 11 test cho \texttt{CourseService} và \\ \texttt{EnrollmentService}, kiểm tra: tính đúng doanh thu và tỷ lệ lấp đầy, ngưỡng \\ \texttt{LowSeatThreshold} theo Options, tìm kiếm, đăng ký thành công thì giảm chỗ, lớp đầy thì từ chối, đăng ký trùng thì từ chối, và \textbf{lỗi giữa transaction thì Rollback toàn bộ}.

            \item \textbf{Layer:} Phần này thuộc tầng kiểm thử, tác động lên \texttt{Service}. Lý do đặt ở Service vì đây là nơi chứa logic nghiệp vụ cần được kiểm chứng độc lập với database.
        \end{itemize}

    \item \textbf{Feature 1 --- Validate Options khi khởi động (ValidateOnStart)}
        \begin{itemize}[label = $\bullet$, leftmargin = 0.3cm]
            \item \textbf{Thêm Data Annotation:} Em bổ sung \texttt{[Range(1, 100)]} cho \texttt{LowSeatThreshold} và \\ \texttt{[Required]} cho \texttt{CenterName} trong \texttt{TrainingCenterConfig}.

            \item \textbf{Kiểm tra ngay khi boot:} Trong \texttt{Program.cs}, em đăng ký Options bằng \\ \texttt{AddOptions<TrainingCenterConfig>().ValidateDataAnnotations().ValidateOnStart()}. Nếu cấu hình sai (ví dụ \texttt{LowSeatThreshold = -5}), ứng dụng \textbf{từ chối khởi động} thay vì chạy với giá trị sai.

            \item \textbf{Layer:} Cấu hình được đặt ở tầng khởi tạo (\texttt{Program.cs}) và Options class, đúng tinh thần Options Pattern ở mức production.
        \end{itemize}

    \item \textbf{Feature 3 --- Concurrency Token chống bán vượt số chỗ (oversell)}
        \begin{itemize}[label = $\bullet$, leftmargin = 0.3cm]
            \item \textbf{Vấn đề thực tế:} Nếu hai người cùng đăng ký chỗ cuối cùng tại cùng thời điểm, cả hai cùng đọc \texttt{CurrentEnrollment} rồi cùng tăng lên, dẫn tới khóa học bị vượt số chỗ tối đa.

            \item \textbf{Giải pháp:} Em thêm cột \texttt{Version} vào Entity \texttt{Course} và cấu hình \\ \texttt{IsConcurrencyToken()} trong \texttt{AppDbContext}. Mỗi lần cập nhật số chỗ, \texttt{Version} cũng tăng. EF Core sinh câu \texttt{UPDATE ... WHERE Version = giá trị cũ}; request commit sau sẽ nhận \texttt{DbUpdateConcurrencyException}, được \texttt{catch} và Rollback với thông báo thân thiện.

            \item \textbf{Migration mới:} Em tạo migration \texttt{AddCourseConcurrencyToken} để thêm cột \texttt{Version} vào database, đồng thời viết một unit test mô phỏng tình huống tranh chấp.

            \item \textbf{Layer:} Rule xử lý tranh chấp đặt ở \texttt{Service} (nghiệp vụ), cấu hình token đặt ở \\ \texttt{DbContext} (Data Layer).
        \end{itemize}

    \item \textbf{Feature 2 --- Trang lọc theo Chuyên ngành và khoảng Học phí}
        \begin{itemize}[label = $\bullet$, leftmargin = 0.3cm]
            \item \textbf{Mục tiêu:} Tạo route \texttt{/Courses/Filter?categoryId=1\&minFee=1000000\&maxFee=4000000} cho phép lọc khóa học theo chuyên ngành (Relationship) và khoảng học phí, đúng yêu cầu Feature 2 của Câu 2.

            \item \textbf{Repository:} Em bổ sung phương thức \texttt{FilterAsync(int? categoryId, decimal? minFee, decimal? maxFee)} vào \texttt{CourseRepository}. Query chạy trực tiếp trên database: \texttt{.Include(c => c.CourseCategory)} để JOIN qua Relationship, \texttt{.Where(...)} lọc theo điều kiện và \texttt{.AsNoTracking()} vì đây là truy vấn chỉ đọc.

            \item \textbf{Service:} \texttt{CourseService.FilterAsync} nhận tham số từ Controller, gọi Repository và map kết quả sang \texttt{CourseListItemViewModel}.

            \item \textbf{Controller:} \texttt{CoursesController.Filter} nhận input từ Query String, gọi Service và trả View --- không chứa logic lọc trực tiếp.

            \item \textbf{View:} Trang \texttt{/Courses/Filter} có form chọn chuyên ngành, nhập học phí tối thiểu và tối đa; kết quả hiển thị dạng card. Dưới kết quả in ra câu EF Core query tương đương để minh họa luồng xử lý.

            \item \textbf{Layer:} Phân chia đúng luồng Controller $\rightarrow$ Service $\rightarrow$ Repository $\rightarrow$ DbContext, mỗi tầng chỉ làm đúng trách nhiệm của mình.
        \end{itemize}

    \item \textbf{Feature Khuyến khích --- Trang Lịch sử đăng ký (\texttt{/Enrollments/History})}
        \begin{itemize}[label = $\bullet$, leftmargin = 0.3cm]
            \item \textbf{Chức năng:} Trang \texttt{/Enrollments/History} liệt kê toàn bộ lịch sử đăng ký, bao gồm tên học viên, tên khóa học, thời gian đăng ký và trạng thái.

            \item \textbf{Relationship:} Truy vấn dùng \texttt{.Include(e => e.Course)} và \texttt{.Include(e => e.Student)} để lấy đầy đủ thông tin qua hai Relationship One-to-Many, cùng với \texttt{AsNoTracking()} vì chỉ đọc.

            \item \textbf{Thống kê:} Hiển thị 3 thẻ thống kê: tổng lượt đăng ký, số học viên duy nhất, số khóa học đã có người đăng ký.
        \end{itemize}

    \item \textbf{Biểu đồ thống kê bằng Chart.js (Tính năng bổ sung)}
        \begin{itemize}[label = $\bullet$, leftmargin = 0.3cm]
            \item \textbf{Trực quan hóa dữ liệu:} Em bổ sung biểu đồ cột doanh thu theo chuyên ngành và biểu đồ tròn tỷ lệ lấp đầy vào trang \texttt{/Courses/Stats}, kèm hiệu ứng count-up cho các con số.

            \item \textbf{Đúng luồng MVC:} Dữ liệu biểu đồ lấy từ \texttt{CourseService.GetStatsAsync()}, đóng gói trong \texttt{CourseStatsViewModel} rồi truyền sang View, không truy vấn trực tiếp trong View.

            \item \textbf{Hỗ trợ sáng/tối:} Màu chữ và lưới của biểu đồ tự đổi theo theme.
        \end{itemize}

    \item \textbf{Hiện đại hóa giao diện --- UI 2025--2026 (Tính năng bổ sung)}
        \begin{itemize}[label = $\bullet$, leftmargin = 0.3cm]
            \item \textbf{Course Card phiên bản mới:} Mỗi card có dải màu nhận diện theo chuyên ngành, thanh tiến trình sĩ số đổi màu theo mức lấp đầy (xanh/vàng/đỏ) và price tag học phí nổi bật.

            \item \textbf{Toast thông báo:} Thay alert chiếm chỗ bằng toast nổi góc phải, áp dụng mô hình PRG (Post-Redirect-Get) sau khi đăng ký thành công.

            \item \textbf{Navbar gọn hơn:} Gom các mục quản trị vào dropdown ``Quản trị'' và đưa ``Thêm mới'' thành nút CTA nổi bật.

            \item \textbf{Chi tiết hiện đại khác:} Design tokens cho theme sáng/tối, View Transitions API, hiệu ứng cuộn xuất hiện cho card, focus ring cho accessibility.
        \end{itemize}
\end{enumerate}

\subsection*{Minh chứng kết quả các tính năng mở rộng}

\textbf{Tính năng 1 --- Unit Test với Fake Repository.} Kết quả chạy \texttt{dotnet test} cho thấy toàn bộ 11 test đều Passed, trong đó có các test kiểm chứng đăng ký thành công, lớp đầy bị từ chối, đăng ký trùng và đặc biệt là Rollback khi lỗi giữa transaction.

\begin{figure}[H]
    \centering
    \includegraphics[width=0.9\textwidth]{img/Lab04/dotnet_test.jpg}
    \caption{Kết quả \texttt{dotnet test}: 11/11 test Passed, chứng minh kiến trúc tách dependency và transaction hoạt động đúng}
\end{figure}

\begin{figure}[H]
    \centering
    \includegraphics[width=0.9\textwidth]{img/Lab04/bonus_1_1.jpg}
    \caption{Chi tiết kết quả test các nghiệp vụ đăng ký (thành công, lớp đầy, đăng ký trùng) --- Test Run Successful}
\end{figure}

\textbf{Tính năng 2 --- Validate Options khi khởi động.} Khi cố tình đặt \texttt{LowSeatThreshold = -5} trong \texttt{appsettings.json}, ứng dụng từ chối khởi động và báo \texttt{OptionsValidationException}. Khi cấu hình hợp lệ, ứng dụng khởi động bình thường.

\begin{figure}[H]
    \centering
    \includegraphics[width=0.9\textwidth]{img/Lab04/II_B1_options_invalid_crash.jpg}
    \caption{Cấu hình sai (\texttt{LowSeatThreshold = -5}): ứng dụng \textbf{từ chối khởi động} với \texttt{OptionsValidationException}}
\end{figure}

\begin{figure}[H]
    \centering
    \includegraphics[width=0.9\textwidth]{img/Lab04/II_B2_options_valid_run.jpg}
    \caption{Cấu hình hợp lệ: ứng dụng khởi động bình thường (so sánh trước/sau)}
\end{figure}

\textbf{Feature 3 --- Concurrency Token.} Migration \texttt{AddCourseConcurrencyToken} đã thêm cột \texttt{Version} vào bảng \texttt{Courses}, và unit test mô phỏng tranh chấp xác nhận request thứ hai bị từ chối, không bán vượt số chỗ.

\begin{figure}[H]
    \centering
    \includegraphics[width=0.9\textwidth]{img/Lab04/II_C1_migrations_list.jpg}
    \caption{Danh sách Migration, trong đó có migration mới \texttt{AddCourseConcurrencyToken} thêm cột \texttt{Version}}
\end{figure}

\begin{figure}[H]
    \centering
    \includegraphics[width=0.9\textwidth]{img/Lab04/sqlite_courses_version_column.jpg}
    \caption{Bảng \texttt{Courses} trong SQLite Viewer --- cột \texttt{Version} đã được thêm vào database bởi migration \texttt{AddCourseConcurrencyToken}}
\end{figure}

\begin{figure}[H]
    \centering
    \includegraphics[width=0.9\textwidth]{img/Lab04/II_C3_concurrency_test_pass.jpg}
    \caption{Unit test \texttt{HaiRequestCungLuc\_RequestSauBiTuChoi\_KhongOversell} Passed: tình huống tranh chấp được xử lý đúng}
\end{figure}

\textbf{Feature 2 --- Lọc theo Chuyên ngành và khoảng Học phí.} Trang \texttt{/Courses/Filter} cho phép lọc khóa học theo chuyên ngành (Relationship) và khoảng học phí. Repository dùng \texttt{Include(c => c.CourseCategory).Where(...).AsNoTracking().OrderBy(...).ToListAsync()} để query trực tiếp database.

\begin{figure}[H]
    \centering
    \includegraphics[width=0.9\textwidth]{img/Lab04/web_courses_filter.jpg}
    \caption{Endpoint \texttt{/Courses/Filter?categoryId=1\&minFee=1000000\&maxFee=4000000}: Lọc khóa học theo chuyên ngành và khoảng học phí, query qua Relationship với AsNoTracking}
\end{figure}

\textbf{Feature Khuyến khích --- Trang Lịch sử đăng ký.} Trang \texttt{/Enrollments/History} liệt kê toàn bộ lịch sử đăng ký, dùng \texttt{Include} qua hai Relationship (\texttt{Course} và \texttt{Student}) với \texttt{AsNoTracking}.

\begin{figure}[H]
    \centering
    \includegraphics[width=0.9\textwidth]{img/Lab04/history_page.jpg}
    \caption{Endpoint \texttt{/Enrollments/History}: Trang lịch sử đăng ký, truy vấn qua hai Relationship (Course và Student) với AsNoTracking}
\end{figure}

\textbf{Tính năng 4 --- Biểu đồ Chart.js.} Trang Thống kê hiển thị biểu đồ doanh thu và tỷ lệ lấp đầy, đồng bộ ở cả chế độ sáng và tối.

\begin{figure}[H]
    \centering
    \includegraphics[width=0.9\textwidth]{img/Lab04/II_D1_stats_light.jpg}
    \caption{Trang Thống kê với biểu đồ Chart.js ở chế độ sáng}
\end{figure}

\begin{figure}[H]
    \centering
    \includegraphics[width=0.9\textwidth]{img/Lab04/II_D2_stats_dark.jpg}
    \caption{Trang Thống kê với biểu đồ Chart.js ở chế độ tối}
\end{figure}

\textbf{Tính năng 5 --- Hiện đại hóa giao diện.} Course Card mới có thanh tiến trình sĩ số và price tag; thông báo đăng ký thành công hiển thị dạng toast; navbar có dropdown quản trị.

\begin{figure}[H]
    \centering
    \includegraphics[width=0.9\textwidth]{img/Lab04/II_D3_course_cards.jpg}
    \caption{Course Card phiên bản mới: dải màu theo chuyên ngành, thanh tiến trình sĩ số và price tag học phí}
\end{figure}

\begin{figure}[H]
    \centering
    \includegraphics[width=0.9\textwidth]{img/Lab04/II_D4_toast_1.jpg}
    \caption{Thông báo đăng ký thành công hiển thị dạng Toast nổi góc phải (mô hình PRG)}
\end{figure}

\begin{figure}[H]
    \centering
    \includegraphics[width=0.9\textwidth]{img/Lab04/II_D5_navbar_dropdown.jpg}
    \caption{Navbar hiện đại với dropdown ``Quản trị'' và nút CTA ``Thêm mới''}
\end{figure}

\section{TRẢ LỜI CÂU HỎI PROBLEM SOLVING (CÂU 3)}

\begin{enumerate}[label = \arabic*.]
    \item \textbf{Nếu đăng ký \texttt{AppDbContext} là Singleton trong web app nhiều request đồng thời, rủi ro gì? Em chọn lifetime nào và vì sao?}

    $\rightarrow$ \textbf{\textit{Trả lời:}}
    \begin{itemize}[leftmargin = 0.6cm, label = -]
        \item \texttt{DbContext} của EF Core \textbf{không thread-safe}. Nếu để Singleton, mọi request dùng chung một instance; khi nhiều request cùng đăng ký khóa học, chúng cùng thao tác trên một \texttt{ChangeTracker} dẫn tới lỗi ``A second operation was started on this context'', dữ liệu bị trộn lẫn và sai lệch.

        \item Trong bài của em, \texttt{AppDbContext} được đăng ký \textbf{Scoped}. Mỗi HTTP request có một \texttt{DbContext} riêng, sống đúng bằng vòng đời request và được giải phóng khi request kết thúc. Đây là lựa chọn chuẩn cho ứng dụng web có Data Layer.
    \end{itemize}

    \item \textbf{Nếu Controller trực tiếp query DbContext và chứa logic nghiệp vụ, việc bảo trì, tái sử dụng và kiểm tra code gặp khó khăn gì?}

    $\rightarrow$ \textbf{\textit{Trả lời:}}
    \begin{itemize}[leftmargin = 0.6cm, label = -]
        \item \textbf{Khó kiểm tra:} Logic nghiệp vụ (như kiểm tra còn chỗ, trừ số chỗ trong transaction) gắn chặt với Controller và database thật, nên không thể viết unit test bằng Fake Repository. Trong bài của em, vì logic nằm ở \texttt{EnrollmentService}, em mới có thể test bằng \\ \texttt{ThrowingEnrollmentRepository}.

        \item \textbf{Khó tái sử dụng:} Nếu sau này có thêm API hoặc màn hình khác cũng cần nghiệp vụ đăng ký, em phải chép lại logic. Đặt ở Service thì chỉ cần gọi lại \texttt{EnrollStudentAsync}.

        \item \textbf{Khó bảo trì:} Controller phình to, lẫn lộn giữa nhận request và xử lý nghiệp vụ, vi phạm nguyên tắc tách trách nhiệm.
    \end{itemize}

    \item \textbf{Rule xác định trạng thái (ngưỡng chỗ trống) nên đặt ở Controller, Service hay Repository?}

    $\rightarrow$ \textbf{\textit{Trả lời:}}
    \begin{itemize}[leftmargin = 0.6cm, label = -]
        \item Em đặt rule ở \textbf{Service}. Trong \texttt{CourseService.GetStatsAsync}, em đọc ngưỡng \\ \texttt{LowSeatThreshold} từ \texttt{IOptions<TrainingCenterConfig>} rồi tính số khóa ``sắp đầy''.

        \item \textbf{Repository} chỉ nên truy vấn dữ liệu thô, không chứa quy tắc nghiệp vụ. \textbf{Controller} chỉ điều phối request. Còn việc ``còn chỗ bao nhiêu thì coi là sắp đầy'' là \textbf{logic nghiệp vụ}, nên thuộc về Service. Nhờ đó khi đổi ngưỡng trong \texttt{appsettings.json}, kết quả thay đổi mà không phải sửa code --- điều này em đã chứng minh bằng unit test \\ \texttt{DoiNguongTrongConfig\_KetQuaThayDoi\_KhongSuaCode}.
    \end{itemize}

    \item \textbf{Nếu một Service tự \texttt{new} Repository bên trong method, vì sao việc thay bằng Fake Repository để kiểm tra logic sẽ khó khăn?}

    $\rightarrow$ \textbf{\textit{Trả lời:}}
    \begin{itemize}[leftmargin = 0.6cm, label = -]
        \item Nếu \texttt{EnrollmentService} tự \texttt{new EnrollmentRepository()} bên trong, nó bị \textbf{phụ thuộc cứng} vào implementation thật và \texttt{DbContext} thật. Khi test, em không có cách nào ``tráo'' Repository khác vào.

        \item Trong bài của em, \texttt{EnrollmentService} nhận \texttt{IEnrollmentRepository} qua constructor (Dependency Injection). Nhờ đó khi test em truyền vào \texttt{ThrowingEnrollmentRepository} để mô phỏng lỗi, hoặc Fake in-memory để kiểm tra logic, mà không cần database. Đây chính là minh chứng cho hình \texttt{ThrowingEnrollmentRepository.cs} ở phần dưới.
    \end{itemize}

    \begin{figure}[H]
        \centering
        \includegraphics[width=0.85\textwidth]{img/Lab04/problem_solving_3.jpg}
        \caption{\texttt{ThrowingEnrollmentRepository} implement \texttt{IEnrollmentRepository}, ném exception ở \texttt{AddAsync} để mô phỏng lỗi giữa transaction (minh chứng cho câu 4 và câu 14)}
    \end{figure}

    \item \textbf{Nếu thêm property mới vào Entity nhưng quên tạo migration, hoặc đã tạo migration nhưng quên chạy database update, ứng dụng có thể gặp lỗi gì?}

    $\rightarrow$ \textbf{\textit{Trả lời:}}
    \begin{itemize}[leftmargin = 0.6cm, label = -]
        \item \textbf{Quên tạo migration:} Model trong code khác với database. Khi truy vấn, EF Core sinh SQL có cột mới (ví dụ \texttt{Version}) nhưng database chưa có cột đó, gây lỗi \texttt{SqliteException: no such column}.

        \item \textbf{Tạo migration nhưng quên update:} Database vẫn ở schema cũ, ứng dụng cũng báo lỗi thiếu cột hoặc bảng. Ngoài ra EF Core có thể báo có migration đang chờ áp dụng (\texttt{pending model changes}).

        \item Vì vậy quy trình đúng là: sửa Entity $\rightarrow$ \texttt{dotnet ef migrations add} $\rightarrow$ \texttt{dotnet ef database update}, như em đã làm với migration \texttt{AddCourseConcurrencyToken}.
    \end{itemize}

    \item \textbf{Nếu trang danh sách đọc 10.000 record nhưng không dùng \texttt{AsNoTracking}, hệ thống tốn thêm tài nguyên thế nào? Khi nào không nên dùng \texttt{AsNoTracking}?}

    $\rightarrow$ \textbf{\textit{Trả lời:}}
    \begin{itemize}[leftmargin = 0.6cm, label = -]
        \item Nếu không dùng \texttt{AsNoTracking}, EF Core tạo \textbf{snapshot theo dõi thay đổi} cho cả 10.000 entity trong \texttt{ChangeTracker}. Điều này tốn nhiều RAM và CPU, làm truy vấn chậm hơn, dù trang danh sách chỉ để hiển thị chứ không sửa gì.

        \item Trong bài của em, các trang chỉ đọc như danh sách khóa học và lịch sử đăng ký đều dùng \texttt{AsNoTracking} để nhẹ và nhanh.

        \item \textbf{Không nên dùng \texttt{AsNoTracking}} khi truy vấn để \textbf{cập nhật} dữ liệu --- ví dụ trong nghiệp vụ đăng ký, em phải truy vấn \texttt{Course} ở chế độ Tracking thì EF Core mới tự sinh \texttt{UPDATE} khi em tăng \texttt{CurrentEnrollment}.
    \end{itemize}

    \item \textbf{Nếu nghiệp vụ đã lưu bản ghi chính nhưng bước giảm số chỗ bị lỗi, dữ liệu sẽ sai thế nào nếu không có transaction?}

    $\rightarrow$ \textbf{\textit{Trả lời:}}
    \begin{itemize}[leftmargin = 0.6cm, label = -]
        \item Không có transaction, bản ghi \texttt{Enrollment} đã được lưu nhưng \texttt{CurrentEnrollment} của khóa học chưa được tăng (hoặc ngược lại). Dữ liệu rơi vào trạng thái \textbf{nửa đúng nửa sai}: số học viên đăng ký thực tế và số chỗ đã dùng không khớp nhau.

        \item Hậu quả là khóa học có thể bị bán vượt số chỗ, hoặc thống kê doanh thu/sĩ số sai. Trong bài của em, vì bọc transaction nên nếu bất kỳ bước nào lỗi, \texttt{RollbackAsync} hủy toàn bộ, đảm bảo dữ liệu luôn nhất quán --- em đã kiểm chứng bằng test \\ \texttt{LoiGiuaChung\_RollbackToanBo}.
    \end{itemize}

    \item \textbf{Hãy mô tả cách em chủ động kiểm tra tình huống rollback. Em tạo dữ liệu đầu vào nào, kiểm tra bảng nào, mong đợi kết quả gì?}

    $\rightarrow$ \textbf{\textit{Trả lời:}}
    \begin{itemize}[leftmargin = 0.6cm, label = -]
        \item \textbf{Dữ liệu đầu vào:} Em tạo một khóa học test có \texttt{CurrentEnrollment = 1}, \texttt{MaxCapacity = 2} (còn chỗ) và một học viên hợp lệ, trên database SQLite in-memory.

        \item \textbf{Tác động lỗi:} Em thay \texttt{IEnrollmentRepository} bằng \\ \texttt{ThrowingEnrollmentRepository} để \texttt{AddAsync} ném exception ngay sau khi đã qua các bước kiểm tra --- mô phỏng lỗi giữa transaction.

        \item \textbf{Kiểm tra bảng:} Em kiểm tra bảng \texttt{Courses} và bảng \texttt{Enrollments}.

        \item \textbf{Kết quả mong đợi:} Hàm trả về thất bại; \texttt{CurrentEnrollment} vẫn là 1 (không bị trừ) và bảng \texttt{Enrollments} có 0 bản ghi. Test \texttt{LoiGiuaChung\_RollbackToanBo} xác nhận đúng các kỳ vọng này. Ngoài ra trên giao diện, khi đăng ký vào lớp đã đầy (PRG-201 30/30), hệ thống cũng hiển thị thông báo từ chối và không thay đổi dữ liệu.
    \end{itemize}

    \begin{figure}[H]
        \centering
        \includegraphics[width=0.85\textwidth]{img/Lab04/problem_solving_2.jpg}
        \caption{Test \texttt{LoiGiuaChung\_RollbackToanBo} theo Arrange - Act - Assert: kiểm tra số chỗ giữ nguyên và 0 enrollment được lưu sau khi lỗi}
    \end{figure}

    \item \textbf{Nếu config được đọc bằng string key ở nhiều nơi, khi đổi tên key hoặc đổi môi trường sẽ có rủi ro gì? Options Pattern giúp giảm rủi ro thế nào?}

    $\rightarrow$ \textbf{\textit{Trả lời:}}
    \begin{itemize}[leftmargin = 0.6cm, label = -]
        \item Đọc bằng string key như \texttt{Configuration["TrainingCenterConfig:LowSeatThreshold"]} rải rác nhiều nơi rất dễ \textbf{gõ sai key} hoặc \textbf{quên sửa} khi đổi tên; lỗi chỉ phát hiện lúc chạy, và giá trị thiếu thường trả về null/0 một cách âm thầm.

        \item \textbf{Options Pattern} gom cấu hình vào class \texttt{TrainingCenterConfig} strongly-typed, truy cập qua property nên được trình biên dịch kiểm tra. Trong bài của em còn thêm \\ \texttt{ValidateDataAnnotations().ValidateOnStart()}, nên nếu cấu hình sai hoặc thiếu, ứng dụng báo lỗi ngay khi khởi động thay vì chạy sai âm thầm.
    \end{itemize}

    \item \textbf{Nếu cùng một request có Service và Repository đều cần DbContext, việc dùng Scoped lifetime giúp ích gì cho tính nhất quán của dữ liệu?}

    $\rightarrow$ \textbf{\textit{Trả lời:}}
    \begin{itemize}[leftmargin = 0.6cm, label = -]
        \item Với Scoped, trong cùng một request, \texttt{EnrollmentService} và \texttt{EnrollmentRepository} được inject \textbf{cùng một instance} \texttt{AppDbContext}.

        \item Nhờ vậy, khi Service mở transaction và Repository thêm \texttt{Enrollment}, cả hai thao tác trên cùng một \texttt{ChangeTracker} và cùng một transaction. Khi gọi \texttt{SaveChangesAsync}, mọi thay đổi (thêm Enrollment, tăng số chỗ) được lưu nhất quán. Nếu mỗi tầng có \texttt{DbContext} riêng, chúng sẽ không chung transaction và dễ gây dữ liệu không đồng bộ.
    \end{itemize}

    \item \textbf{Seed Data có thể gây trùng dữ liệu trong trường hợp nào? Em thiết kế seed thế nào để không tạo dữ liệu lặp?}

    $\rightarrow$ \textbf{\textit{Trả lời:}}
    \begin{itemize}[leftmargin = 0.6cm, label = -]
        \item Seed có thể gây trùng nếu mỗi lần khởi động lại chèn dữ liệu mới mà không kiểm tra đã tồn tại, hoặc gán \texttt{Id} không cố định khiến EF coi là bản ghi mới.

        \item Trong bài của em, em dùng \texttt{modelBuilder.Entity<T>().HasData(...)} với \texttt{Id} cố định (ví dụ \texttt{Id = 1, 2, 3, 4}). Cơ chế \texttt{HasData} gắn với Migration nên dữ liệu seed được quản lý theo khóa chính: chạy update nhiều lần cũng không tạo bản ghi lặp, vì EF so khớp theo \texttt{Id}.
    \end{itemize}

    \item \textbf{Nếu mở rộng transaction từ một item thành nhiều item, cần kiểm tra điều kiện nào trước khi commit để tránh giao dịch nửa đúng nửa sai?}

    $\rightarrow$ \textbf{\textit{Trả lời:}}
    \begin{itemize}[leftmargin = 0.6cm, label = -]
        \item Nếu một lần đăng ký nhiều khóa học, trước khi commit em cần kiểm tra \textbf{tất cả} các điều kiện cho \textbf{từng khóa}: mỗi khóa còn chỗ, học viên chưa đăng ký trùng khóa đó, và mỗi khóa không bị tranh chấp (concurrency token).

        \item Chỉ khi \textbf{toàn bộ} các khóa đều hợp lệ mới \texttt{CommitAsync}. Nếu chỉ một khóa không đủ chỗ, phải \texttt{RollbackAsync} toàn bộ để không xảy ra trường hợp đăng ký được khóa này nhưng khóa kia thất bại --- đảm bảo tính ``tất cả hoặc không gì cả'' của transaction.
    \end{itemize}

    \item \textbf{Hãy đề xuất cách chia Controller, Service, Repository cho chức năng lọc dữ liệu theo Relationship và khoảng giá trị (Câu 2).}

    $\rightarrow$ \textbf{\textit{Trả lời:}}
    \begin{itemize}[leftmargin = 0.6cm, label = -]
        \item \textbf{Controller:} chỉ nhận input từ Query String (ví dụ \texttt{categoryId}, \texttt{keyword}), gọi Service và trả View. Không chứa logic lọc.

        \item \textbf{Service:} nhận tham số, gọi Repository lấy dữ liệu, xử lý logic (lọc theo chuyên ngành, theo khoảng giá trị) và map sang ViewModel.

        \item \textbf{Repository:} chịu trách nhiệm query database, dùng \texttt{Include} cho Relationship và \\ \texttt{AsNoTracking} vì đây là truy vấn chỉ đọc.

        \item Cách chia này giữ đúng luồng \texttt{Controller $\rightarrow$ Service $\rightarrow$ Repository $\rightarrow$ DbContext} và giúp mỗi tầng dễ test, dễ bảo trì.
    \end{itemize}

    \item \textbf{Hãy mô tả một kịch bản kiểm tra Service bằng Fake Repository theo ba bước Arrange - Act - Assert.}

    $\rightarrow$ \textbf{\textit{Trả lời:}} Em lấy ví dụ test \texttt{EnrollStudentAsync\_LoiGiuaChung\_RollbackToanBo} trong bài của em:
    \begin{itemize}[leftmargin = 0.6cm, label = -]
        \item \textbf{Arrange:} Chuẩn bị database SQLite in-memory với một khóa học còn chỗ và một học viên; khởi tạo \texttt{EnrollmentService} với \texttt{ThrowingEnrollmentRepository} (Fake ném exception ở \texttt{AddAsync}).

        \item \textbf{Act:} Gọi \texttt{service.EnrollStudentAsync(courseId, studentId)} --- nghiệp vụ sẽ lỗi giữa transaction.

        \item \textbf{Assert:} Kiểm tra kết quả trả về là thất bại; \texttt{CurrentEnrollment} của khóa học vẫn giữ nguyên (không bị trừ) và bảng \texttt{Enrollments} có 0 bản ghi --- chứng minh transaction đã Rollback toàn bộ.
    \end{itemize}
\end{enumerate}
